\chapter{Comparación de modelos y estrategia de trabajo}
\label{cap:comparacion-modelos}

El disco y el semiplano representan la misma geometría, pero hacen visibles aspectos distintos. Elegir coordenadas convenientes suele simplificar una demostración o un cálculo.

\section{Qué modelo conviene usar}

\begin{table}[htbp]
  \centering
  \caption{Dos modelos del plano hiperbólico.}
  \label{tab:modelos}
  \small
  \begin{tabular}{p{0.18\textwidth}p{0.34\textwidth}p{0.34\textwidth}}
    \toprule
    Propiedad & Disco de Poincaré & Semiplano superior \\
    \midrule
    Dominio & \( |z|<1 \) & \( y>0 \) \\
    Frontera ideal & \( |z|=1 \) & eje real y \(\infty\) \\
    Geodésicas & diámetros y arcos ortogonales & verticales y semicírculos \\
    Ventaja visual & simetría circular & transformaciones matriciales \\
    \bottomrule
  \end{tabular}
\end{table}

El disco resulta cómodo para visualizar configuraciones globales y polígonos ideales. El semiplano es útil para integrar a lo largo de verticales y para aplicar transformaciones de \(\mathrm{PSL}(2,\mathbb R)\). La tabla \cref{tab:modelos} resume las diferencias de coordenadas, no una diferencia geométrica.

\section{Un método para resolver problemas}

Al abordar un ejercicio de geometría hiperbólica, puede servir el siguiente orden:
\begin{enumerate}
  \item Identifica si la pregunta trata de longitud, ángulo, área, incidencia o curvatura.
  \item Escoge el modelo que haga más sencilla la figura o la expresión.
  \item Usa una isometría para mover puntos a posiciones simples.
  \item Comprueba qué fórmula métrica corresponde al modelo.
  \item Escribe la cadena de igualdades y señala las hipótesis.
  \item Interpreta el resultado en términos geométricos y verifica que tenga sentido en casos límite.
\end{enumerate}

\section{Cómo seguir leyendo}

Los capítulos anteriores desarrollan modelos, geodésicas, áreas y teselaciones con ejemplos elementales. Para continuar, estas obras ofrecen una ruta de lectura:

\begin{itemize}
  \item Anderson presenta la geometría hiperbólica en un texto introductorio dedicado al tema \cite{anderson2005}.
  \item Greenberg desarrolla el contexto histórico y axiomático de las geometrías no euclidianas \cite{greenberg2008}.
  \item Stillwell estudia superficies y teselaciones \cite{stillwell1992}.
  \item Beardon y Ratcliffe ofrecen referencias más avanzadas sobre grupos discretos y variedades hiperbólicas \cite{beardon1983,ratcliffe2006}.
\end{itemize}

\begin{remark}
La bibliografía es un punto de partida, no una lista exhaustiva. Al ampliar estas notas conviene anotar qué resultado se tomó de cada obra y añadir la cita junto al argumento correspondiente.
\end{remark}
